% §14.5 Construction Time and Trade-offs
%
% Cross-reference resolution status:
%
% DEFINED here:
%   \label{sec:construction-time}                -- this section
%
% ALREADY DEFINED (chapter 5):
%   \Cref{def:bloom-filter}                      -- chapter 5 §5.2
%   \Cref{alg:bloom-insert-query}                -- chapter 5 §5.2 pseudocode
%
% ALREADY DEFINED (chapter 8):
%   \Cref{prop:bloom-overhead}                   -- chapter 8 §8.4
%
% ALREADY DEFINED (chapter 14, earlier sections):
%   \Cref{def:bhf}                               -- §14.2
%   \Cref{alg:bhf-insert-query}                  -- §14.2 pseudocode
%   \Cref{sec:bhf-construction}                  -- §14.2
%   \Cref{sec:space-analysis}                    -- §14.3
%   \Cref{thm:bhf-optimality}                    -- §14.3
%   \Cref{prop:bhf-fpr}                          -- §14.4
%
% FORWARD REFS within chapter 14:
%   \Cref{sec:ch14-bridge}                       -- §14.6 (Task 7)
%   \Cref{sec:bib-notes-ch14}                    -- §14.7 (Task 8)
%
% §14.5 covers the remaining performance dimensions: build time
% (perfect-hash construction is O(n) expected, fingerprint table is
% O(n) one-pass); query time O(1); trade-offs versus Bloom (static
% set, no in-place insertions); decision tree for when to use which.
% Perfect-hash algorithm names appear in prose; citations are deferred
% to §14.7 bibliographic notes.

\section{Construction Time and Trade-offs}
\label{sec:construction-time}

The space analysis of \Cref{sec:space-analysis} and the FPR analysis
of \Cref{sec:fpr-analysis} settled the BHF's accuracy and bit
budget. The remaining performance dimensions are build time and
query latency, and the structural property that distinguishes the
BHF from the Bloom filter as a deployable data structure: the BHF
is static. Inserting a new element after build requires recomputing
the perfect hash, which is the dominant build-time cost. This
section quantifies the costs, contrasts them with the Bloom
filter's, and offers a decision rule for choosing between the two
constructions in practice.

\subsection*{Building the perfect hash}

The perfect hash $\phi : A \to [n]$ is constructed from $A$ at build
time. Several algorithm families achieve expected $O(n)$
construction time, with different leading constants and different
output-size trade-offs. The classical Fredman-Komlos-Szemeredi
construction (1984) achieves $O(n)$ expected time with $O(n \log n)$
output bits. The Hagerup-Tholey scheme (2001) reduces the output to
$O(n)$ bits. The modern displacement-based family, with MWHC and
BBHash as the best-known representatives, achieves $O(n)$ expected
time with $O(n)$ output bits at small leading constants (between
$2$ and $4$ bits per element). The asymptotic family of
Belazzougui, Botelho, and Dietzfelbinger (2009) achieves $O(n)$
expected time with $O(n \log \log \log n / \log n)$ output bits, the
sublinear bound that \Cref{sec:space-analysis} used. The algorithm
names are sketched here; the bibliographic entries appear in
\Cref{sec:bib-notes-ch14}.

The expected-time guarantee comes from a randomization step: each
algorithm samples a random parameter (a seed, a displacement vector,
a partitioning hash) and retries on failure. The expected number of
retries is constant, so the expected total construction time is
linear in $n$. The worst case can be longer but is vanishingly rare
under standard assumptions on the input distribution. For
pedagogical purposes the BHF treats the perfect-hash construction
as an expected-$O(n)$ subroutine that returns a representation
satisfying the requirements of \Cref{def:bhf}; the choice of which
specific algorithm to use affects constants but not the asymptotic
behaviour of the BHF as a whole.

The constant for the perfect-hash construction matters in practice
because it dominates BHF build time. The fingerprint-table fill, in
contrast, is a single pass over $A$ at one hash evaluation per
element and is fast (memory-bandwidth-bound, not compute-bound).
Modern displacement-based implementations achieve sustained
construction throughputs on the order of $10$ to $50$ million
elements per second on commodity hardware. The BHF's build time is
therefore on the order of seconds for million-element sets and on
the order of minutes for billion-element sets, dominated by the
perfect-hash construction throughout that range. Both the BHF and
the Bloom filter scale linearly in $n$; the difference is that the
Bloom filter's per-element constant is much smaller (a few hash
evaluations and bit-writes, no randomized retries) so its absolute
build time is lower at the same $n$. The BHF pays the perfect-hash
construction cost up front in exchange for the per-element space
savings the chapter has documented.

\subsection*{Building the fingerprint table}

Once $\phi$ is in hand, the fingerprint table $F$ is built in one
pass over $A$:
\[
  \text{for each } a \in A: \quad F[\phi(a)] \;\coloneqq\; g(a),
\]
the construction pass already named in \Cref{sec:bhf-construction}.
The pass evaluates $\phi(a)$ once per element (a constant-time
lookup once $\phi$ has been constructed), evaluates $g(a)$ once per
element (one hash evaluation), and writes one $r$-bit slot per
element. Total work is $O(n)$ time and $O(n \cdot r)$ output bits,
matching the fingerprint-storage bound of \Cref{sec:space-analysis}.

The fingerprint-table fill is embarrassingly parallel: each
iteration writes a distinct slot (because $\phi$ is a bijection on
$A$), so threads can split $A$ into chunks and fill the table
without coordination. The perfect-hash construction is less easily
parallelized but admits parallel variants in the published
implementations. The fingerprint-fill stage and the perfect-hash
stage together compose into a build pipeline whose end-to-end
throughput on commodity hardware is in the same range as the
perfect-hash construction alone, since the fingerprint fill is
memory-bound and overlaps well with the perfect-hash output
streaming.

The order in which the two stages run matters: the perfect hash
must be available before the fingerprint table can be filled,
because each fingerprint write addresses slot $\phi(a)$. The build
pipeline is therefore sequential at the stage level (perfect-hash
first, fingerprints second) but parallel within each stage. A
streaming variant that interleaves the two stages is possible if
the perfect-hash algorithm supports incremental output, but the
displacement-based families typically emit the perfect-hash
representation in batches rather than per-element, so the
interleaving offers limited speedup. For most practical builds the
two-stage pipeline is the right organization.

\subsection*{Query time}

The query algorithm of \Cref{alg:bhf-insert-query} runs in $O(1)$
time. It evaluates the perfect hash $\phi(x)$ once (one
constant-time data-structure lookup), evaluates the fingerprint
hash $g(x)$ once (one hash evaluation), reads one $r$-bit slot of
the table (one memory access), and compares two $r$-bit strings
(constant time when $r$ fits in a machine word, which it does for
all practical FPR targets). The total work is constant in $n$ and
constant in the universe size, dominated by the two hash
evaluations and the one memory read.

Compared to the Bloom filter's query cost of $k$ hash evaluations
and $k$ memory reads (with $k = (m/n) \ln 2$ at the optimum, which
is $k \approx \log_2(1/\fprate) / \log_2 e \cdot \ln 2 = \ln 2 \cdot
\log_2(1/\fprate) \approx 0.69 \log_2(1/\fprate)$), the BHF's
query cost is asymptotically lower by a factor of roughly $k/2$.
At $\fprate = 0.01$, $k \approx 4.6$, so the BHF performs roughly
$2.3$ times fewer hash evaluations per query and roughly the same
factor fewer memory accesses. At tighter targets the gap widens:
at $\fprate = 10^{-6}$, $k \approx 14$, so the BHF performs
roughly $7$ times fewer hash evaluations per query. The BHF is
denser per bit (the chapter's main result) and faster per query
(this section's secondary result); the price for both is the
static-set restriction the next subsection addresses.

\subsection*{Trade-offs versus the Bloom filter}

The chapter's two structural results so far place the BHF strictly
above the Bloom filter on the two axes of accuracy and space at the
same target: the BHF achieves the entropy floor of
\Cref{thm:bhf-optimality} and beats the Bloom filter's $\log_2 e$
overhead of \Cref{prop:bloom-overhead}; the BHF satisfies the FPR
target of \Cref{prop:bhf-fpr} under the same uniformity assumption
the Bloom filter requires. On the third axis, the cost-of-update,
the Bloom filter is strictly above the BHF. The Bloom filter
supports streaming insertions in $O(k)$ time per insertion, with
no rebuilding required. The BHF requires rebuilding the perfect
hash whenever the designated set $A$ changes by addition, and the
rebuild cost is the $O(n)$ perfect-hash construction. The asymmetry
is the source of the trade-off.

Two operations need separate accounting. \emph{Insertion} of a new
element $a' \notin A$ into a Bloom filter is $O(k)$: evaluate $k$
hash functions and write $k$ bits. Insertion into a BHF requires
rebuilding the perfect hash $\phi$ on $A \cup \{a'\}$ and refilling
the fingerprint table, total $O(n)$. \emph{Deletion} of an element
$a \in A$ from a Bloom filter is not supported by the standard
construction (deletion would require unsetting bits that other
elements may also use). Deletion from a BHF requires rebuilding the
perfect hash on $A \setminus \{a\}$, also $O(n)$. Neither structure
supports cheap deletions; the BHF supports expensive deletions,
the standard Bloom filter supports none at all.

\Cref{tab:bhf-bloom-tradeoffs} summarizes the trade-offs.

\begin{table}[h]
\centering
\caption{BHF versus standard Bloom filter trade-offs at the same
target $\fprate$ and the same $n$. Space and query times exclude
the sublinear perfect-hash overhead the BHF inherits.}
\label{tab:bhf-bloom-tradeoffs}
\begin{tabular}{l|cc}
\hline
                                 & BHF                          & Bloom              \\
\hline
  Bits per element               & $\lceil -\log_2 \fprate \rceil$ & $\log_2(1/\fprate) \cdot \log_2 e$ \\
  Build time (whole set)         & $O(n)$ expected              & $O(k n) = O(n \log_2(1/\fprate))$ \\
  Insertion (post-build)         & $O(n)$ (full rebuild)        & $O(k) = O(\log_2(1/\fprate))$ \\
  Deletion                       & $O(n)$ (full rebuild)        & not supported       \\
  Query                          & $O(1)$ ($2$ hashes, $1$ read)& $O(k)$ ($k$ hashes, $k$ reads) \\
  False-negative rate            & $0$                          & $0$                 \\
  False-positive rate            & $\le 2^{-r} \le \fprate$     & $\fprate$ (asymptotic) \\
\hline
\end{tabular}
\end{table}

The table separates the BHF's structural advantages (denser, faster
queries, supports deletion via rebuild) from its structural cost
(insertions and deletions are $O(n)$ rather than $O(k)$). The
trade-off is between batch and stream: the BHF is the right choice
when the designated set is built once and queried many times; the
Bloom filter is the right choice when the set evolves and the
insertion rate is high relative to the query rate. For applications
where the set is rebuilt periodically (daily, hourly), the BHF's
$O(n)$ rebuild cost is amortized over the inter-rebuild query
volume and the BHF still wins on aggregate. For applications where
insertions arrive interleaved with queries at comparable rates, the
Bloom filter's $O(k)$ insertion cost is decisive and the BHF cannot
compete without architectural changes (e.g., the dynamic variant
the exercises sketch, which buffers insertions and rebuilds in
chunks).

\subsection*{Decision rule}

The choice between BHF and Bloom resolves on two questions: is the
set static or streaming, and is space the binding constraint? Three
combinations of answers cover the practical cases.

\smallskip
\noindent\emph{Static set, space-critical: BHF.} When the designated
set is fixed at build time (or rebuilt on a schedule), and the
deployment is constrained by memory or storage budget (the data
structure must fit in RAM on the smallest available machine, or
must be transmitted over a slow network), the BHF is the right
choice. The space savings of $\log_2 e \approx 1.44$ versus the
Bloom filter translate directly to a smaller memory footprint or a
shorter transmission time. Examples: a network protocol that ships
membership filters to clients on connection establishment; a
read-only database index loaded into RAM at startup and queried
millions of times before reload; a content-delivery network's URL
cache populated nightly and queried throughout the day.

\smallskip
\noindent\emph{Streaming set, space-tolerant: Bloom.} When elements
arrive continuously and the data structure must absorb them in
real time, and when memory is plentiful relative to the working
set, the Bloom filter is the right choice. The $44\%$ space
overhead is a small price for $O(k)$ amortized insertion versus
$O(n)$ rebuild. Examples: a deduplication filter for streaming
ingest pipelines; an in-memory cache of recently-seen request IDs;
a probabilistic check in a routing layer that must handle line-rate
traffic.

\smallskip
\noindent\emph{Hybrid: tiered or dynamic.} When the set has both a
large static portion and a small streaming portion (the common case
in practice), a tiered structure combines the two: a BHF for the
static base layer and a small Bloom filter for the streaming delta,
with periodic merges that absorb the delta into a rebuilt BHF. The
hybrid achieves most of the BHF's space efficiency on the static
portion while accommodating streaming updates through the delta
layer. The merge cost is the BHF rebuild on $|A| + |\Delta|$
elements, amortized over the inter-merge interval. The exercises
return to the tiered organization; for the present chapter the
two pure cases are the operating decision rule, and the hybrid is
a refinement layered on top.

The decision rule connects the chapter's space-optimality result
to the deployment question. The BHF is not a universal replacement
for the Bloom filter; it is the optimal choice when the set is
static and the bit budget matters. The chapter's headline result
(BHF closes the $\log_2 e$ gap that the Bloom filter pays) is
delivered in exchange for the static-set restriction. The next
section bridges from the chapter's construction to its successor,
the BPHF of chapter~15, which exchanges the FPR-zero half of the
perfect-filter line for the FNR-zero side and answers the
ranked-search dual of the chapter's membership question.
